\documentclass[11pt]{article}
\usepackage{fontspec}
\setmainfont{Latin Modern Roman}
\usepackage{amsmath,amssymb,mathtools,unicode-math}
\setmathfont{Latin Modern Math}
\usepackage{xurl,hyperref}
\hypersetup{hidelinks,pdftitle={Problem statement}}
\newcommand{\sref}[2]{\hyperlink{source:#1:#2}{[S#2]}}
\newcommand{\cataloguescope}[1]{\paragraph{Mathematical question.}#1}
\begin{document}
\section{Open Problems in Integrable Systems — Bi-Poisson vector spaces}\label{statement}

\subsection{Definitions and mathematical statement}

\paragraph{Shared notation.}$\mathbb N=\{1,2,\ldots\}$, and $\mathbb Z,\mathbb Q,\mathbb R,\mathbb C$ denote the integers, rationals, reals and complex numbers. Any other conventions are specified locally.


Let $V$ be a finite-dimensional complex vector space with alternating bilinear forms $A,B$, and let $J=\{A+\lambda B:\lambda\in\mathbb C\}\cup\{B\}$. Its generic corank $c$ is the minimum corank of these forms. Define \[\operatorname{Aut}(V,J)=\{T\in GL(V):A(Tu,Tv)=A(u,v),\ B(Tu,Tv)=B(u,v)\ \forall u,v\in V\}.\] Thus the automorphisms preserve each form, not merely the pencil as an unparameterized set. The Grassmannian $\operatorname{Gr}_k(V)$ consists of $k$-dimensional linear subspaces.
\paragraph{Statement recorded in the supplied source.}
Describe the orbits of $\operatorname{Aut}(V,J)$ on $V$ and, for every $0\le k\le\dim V$, on $\operatorname{Gr}_k(V)$. Determine orbit invariants and the fixed points of the Grassmannian action, that is, the subspaces invariant under every pencil automorphism. \sref{statement}{2}
\subsection{Short English statement}
Classify vectors and subspaces under the linear group preserving a pair of alternating forms, including all invariant subspaces.
\subsection{Sources}
\begin{description}
\item[\hypertarget{source:statement:1}{S1}] ULAM.AI and the UnsolvedMath Contributors, UnsolvedMath: A Curated Collection of Open Mathematics Problems, catalogue entry. CC BY 4.0. \url{https://huggingface.co/datasets/ulamai/UnsolvedMath}.
\item[\hypertarget{source:statement:2}{S2}] A. Bolsinov, V. S. Matveev, E. Miranda and S. Tabachnikov, Open Problems, Questions, and Challenges in Finite-Dimensional Integrable Systems (2018), arXiv:1804.03737, Problem 5.2 and automorphism definition. \url{https://arxiv.org/abs/1804.03737}.
\end{description}
\label{end:statement}
\end{document}
